\documentclass[12pt]{article}

\usepackage{amsmath, amssymb}
\usepackage{geometry}
\usepackage{setspace}
\usepackage{titlesec}

\geometry{margin=1in}
\setstretch{1.2}

\title{Lecture Scribe\\
\large CSE400 -- Fundamentals of Probability in Computing\\
Lecture 5: Bayes’ Theorem, Random Variables, and Probability Mass Function}
\author{Instructor: Dhaval Patel, PhD}
\date{January 20, 2026}

\begin{document}

\maketitle

\section{Introduction / Scope}

Topics covered in the lecture:

\begin{itemize}
\item Bayes’ Theorem
    \begin{itemize}
    \item Weighted Average of Conditional Probabilities
    \item Learning by Example
    \item Law of Total Probability and Bayes Formula
    \end{itemize}
\item Random Variables
    \begin{itemize}
    \item Motivation and Concept
    \item Examples
    \end{itemize}
\item Probability Mass Function (PMF)
    \begin{itemize}
    \item Concept
    \item Examples
    \end{itemize}
\end{itemize}

\section{Bayes’ Theorem}

\subsection{Weighted Average of Conditional Probabilities}

\subsubsection*{Decomposition of an Event}

Let $A$ and $B$ be events.

Event $A$ may be expressed as:

\[
A = AB \cup AB^c
\]

Reason: For an outcome to be in $A$, it must either:
\begin{itemize}
\item be in both $A$ and $B$, or
\item be in $A$ but not in $B$.
\end{itemize}

Since $AB$ and $AB^c$ are mutually exclusive, by Axiom 3:

\[
Pr(A) = Pr(AB) + Pr(AB^c)
\]

Using conditional probability:

\[
Pr(A) = Pr(A|B)Pr(B) + Pr(A|B^c)[1 - Pr(B)]
\]

Statement from lecture:  
The probability of event $A$ is a weighted average of the conditional probabilities, with weights equal to the probability of the conditioning event.

\subsection{Learning by Example}

\subsubsection*{Example 3.1 (Part 1)}

\textbf{Problem}

Population classification:
\begin{itemize}
\item Accident prone
\item Not accident prone
\end{itemize}

Given:

\[
Pr(\text{accident in 1 year} \mid \text{accident prone}) = 0.4
\]

\[
Pr(\text{accident in 1 year} \mid \text{not accident prone}) = 0.2
\]

\[
Pr(\text{accident prone}) = 0.3
\]

Find probability that a new policyholder has an accident within a year.

\subsubsection*{Solution}

Let:
\begin{itemize}
\item $A_1$: policyholder has an accident within a year
\item $A$: policyholder is accident prone
\end{itemize}

\[
Pr(A_1) = Pr(A_1|A)Pr(A) + Pr(A_1|A^c)Pr(A^c)
\]

Substitute values:

\[
Pr(A_1) = (0.4)(0.3) + (0.2)(0.7)
\]

\[
Pr(A_1) = 0.26
\]

\subsubsection*{Example 3.1 (Part 2)}

\textbf{Problem}

Given a policyholder had an accident within a year, find probability that the person is accident prone.

\subsubsection*{Solution}

\[
Pr(A|A_1) = \frac{Pr(AA_1)}{Pr(A_1)}
\]

Using:

\[
Pr(AA_1) = Pr(A)Pr(A_1|A)
\]

\[
Pr(A|A_1) =
\frac{Pr(A)Pr(A_1|A)}{Pr(A_1)}
\]

Substitute values:

\[
Pr(A|A_1) = \frac{(0.3)(0.4)}{0.26}
\]

\[
= \frac{6}{13}
\]

\subsection{Law of Total Probability and Bayes Formula}

Statement from lecture:

\[
Pr(AB_i) = Pr(B_i|A)Pr(A)
\]

Using this relation leads to \textbf{Bayes Formula (Proposition 3.1)}.

Terminology:
\begin{itemize}
\item $Pr(B_i)$: \textbf{apriori probability}
\item $Pr(B_i|A)$: \textbf{posteriori probability} (after observing event $A$)
\end{itemize}

\subsection{Example 3.2 (Cards Problem)}

\subsubsection*{Problem}

Three cards:
\begin{enumerate}
\item Both sides red (RR)
\item Both sides black (BB)
\item One red, one black (RB)
\end{enumerate}

Procedure:
\begin{itemize}
\item One card selected randomly
\item Card placed on ground
\item Upper side is red
\end{itemize}

Find probability that the other side is black.

\subsubsection*{Solution}

Define events:
\begin{itemize}
\item RR: chosen card is all red
\item BB: chosen card is all black
\item RB: chosen card is red--black
\item R: upturned side is red
\end{itemize}

Desired probability:

\[
Pr(RB|R) =
\frac{Pr(RB \cap R)}{Pr(R)}
\]

Using Bayes form:

\[
Pr(RB|R)=
\frac{Pr(R|RB)Pr(RB)}
{Pr(R|RR)Pr(RR)+Pr(R|RB)Pr(RB)+Pr(R|BB)Pr(BB)}
\]

Substitute values:

\[
\frac{(1/2)(1/3)}
{(1)(1/3)+(1/2)(1/3)+(0)(1/3)}
\]

Result:

\[
= \frac{1}{3}
\]

\subsubsection*{Enumeration Explanation}

Possible equally likely outcomes:

\[
R_1, R_2, B_1, B_2, R_3, B_3
\]

Other side is black only when outcome is $R_3$.

Conditional probability:

\[
Pr(R_3 \mid R_1 \text{ or } R_2 \text{ or } R_3) = \frac{1}{3}
\]

\section{Random Variables}

\subsection{Motivation and Concept}

Definition:

A \textbf{random variable} is a real-valued function defined on the sample space.  
Values are determined by outcomes of an experiment.  
Probabilities are assigned to possible values of random variables.

\subsubsection*{Examples from Lecture}

\textbf{Experiment 1:} Dice tossing — interest in the sum of the two dice.

\textbf{Experiment 2:} Coin flipping — interest in total number of heads.

\subsubsection*{Distribution Visualization}

Distribution of a random variable can be represented as a bar diagram:
\begin{itemize}
\item x-axis: values of random variable
\item bar height at value $a$: $Pr[X=a]$
\item Probabilities computed from corresponding events in sample space
\end{itemize}

\subsection{Example: Tossing 3 Fair Coins}

Let $Y$ = number of heads.

Possible values:

\[
Y \in \{0,1,2,3\}
\]

Probabilities:

\[
P\{Y=0\} = P\{(t,t,t)\} = \frac{1}{8}
\]

\[
P\{Y=1\} = P\{(t,t,h),(t,h,t),(h,t,t)\} = \frac{3}{8}
\]

\[
P\{Y=2\} = P\{(t,h,h),(h,t,h),(h,h,t)\} = \frac{3}{8}
\]

\[
P\{Y=3\} = P\{(h,h,h)\} = \frac{1}{8}
\]

Since $Y$ must take one of these values, the probabilities sum to 1.

\section{Probability Mass Function (PMF)}

\subsection{Definition}

A random variable that can take at most a countable number of values is called \textbf{discrete}.

Let $X$ be a discrete random variable with range:

\[
R_X = x_1, x_2, x_3, \ldots
\]

The function:

\[
p(x_k) = Pr(X=x_k)
\]

is called the \textbf{Probability Mass Function (PMF)} of $X$.

Since $X$ must take one of the values $x_k$:

\[
\sum_k p(x_k) = 1
\]

\subsection{Example}

Given PMF:

\[
p(i)=c\lambda^i,\quad i=0,1,2,\ldots
\]

Find:
\begin{itemize}
\item $P\{X=0\}$
\item $P\{X>2\}$
\end{itemize}

(The lecture presents the problem statement; intermediate solution steps are not shown in the slides and are therefore not expanded here.)

\section{End of Lecture}

Class participation instruction:  
Switch to Campuswire for discussion and participation.

\end{document}
